% TOP_PROBLEM: {"id": 5500068, "title": "Rolling a Die over a Labeled Board", "category_id": 2, "category": {"id": 2, "name": "combinatorics", "display_name": "Combinatorics", "description": "Counting problems, graph theory, discrete structures.", "slug": "combinatorics", "order_index": 2, "created_at": "2026-07-31T15:26:25.671Z"}, "status": "open"}
% FIRST_POSTED: null
% ENTRY_KIND: statement_only
% Imported from: batch07.tex; UnsolvedMath ID 5500068
\documentclass[11pt]{article}
\usepackage{amsmath,amssymb,mathtools}
\usepackage{fontspec,unicode-math}
\setmainfont{Latin Modern Roman}
\setsansfont{Latin Modern Sans}
\setmathfont{Latin Modern Math}
\usepackage{xurl,hyperref}
\newcommand{\sref}[2]{\hyperlink{source:#1:#2}{[S#2]}}
\newcommand{\cataloguescope}[1]{\paragraph{Mathematical question.}#1}
\begin{document}
% UnsolvedMath ID 5500068; source catalogue code AMR-054-0068
\setcounter{section}{5500067}
\section{Rolling a Die over a Labeled Board}\label{5500068}
{\small\sffamily UnsolvedMath ID 5500068 · Combinatorics}
\subsection{Definitions and mathematical statement}

\paragraph{Shared notation.}$\mathbb N=\{1,2,\ldots\}$, and $\mathbb Z,\mathbb Q,\mathbb R,\mathbb C$ denote the integers, rationals, reals and complex numbers. Any other conventions are specified locally.


A unit cube has its faces labeled $1,\ldots,6$ as a die. It starts on a square of the integer grid at the corner of a rectangular board $R$, with edges parallel to the axes. Every square of $R$ has a prescribed label in $\{1,\ldots,6\}$. A move rolls the cube by a quarter-turn over one of its bottom edges to an adjacent board square. There are no free cells: each board square must be visited exactly once, and its label must equal the cube\textquotesingle s upper face on that visit.
\paragraph{Statement recorded in the supplied source.}
Determine the computational complexity of deciding whether a legal rolling sequence exists for a fully labeled rectangular board with the prescribed starting die. If a solution exists, is the fully labeled board uniquely rollable? \sref{5500068}{2}
\subsection{Short English statement}
Determine the computational complexity of deciding whether a legal rolling sequence exists for a fully labeled rectangular board with the prescribed starting die. If a solution exists, is the fully labeled board uniquely rollable?
\subsection{Sources}
\begin{description}
\item[\hypertarget{source:5500068:1}{S1}] ULAM.AI and the UnsolvedMath Contributors, UnsolvedMath: A Curated Collection of Open Mathematics Problems, record 5500068 (AMR-054-0068). CC BY 4.0. \url{https://huggingface.co/datasets/ulamai/UnsolvedMath}.
\item[\hypertarget{source:5500068:2}{S2}] Erik D. Demaine, Joseph S. B. Mitchell and Joseph O\textquotesingle Rourke (editors), \emph{The Open Problems Project}, maintained original source, Problem 68, Statement and complete Partial/Related Results. \url{https://github.com/edemaine/topp/blob/main/Problems/P.000068}.
\end{description}
\label{end:5500068}
\end{document}
